\chapter{Discusión}

\section{Cumplimiento del problema planteado}

La solución resuelve el problema principal: conectar un catálogo técnico en OpenMetadata con una salida interoperable DCAT-AP-ES validable. La plataforma no se limita a documentar un mapeo, sino que implementa un flujo ejecutable que descubre activos técnicos, permite curación funcional, aplica metadatos, exporta JSON-LD y valida el resultado. El Anexo~\ref{ap:demo-operativa} narra ese flujo de extremo a extremo desde la consola web, hasta dejar el archivo RDF listo para entregarlo a un harvester compatible.

La decisión de usar PostgreSQL de referencia en lugar de cosechar CKAN refuerza la defendibilidad del trabajo. CKAN habría permitido transformar metadatos ya publicados; PostgreSQL y OpenMetadata permiten demostrar gobierno desde activos técnicos, que es más coherente con las competencias de ingeniería, Big Data y cloud del máster.

\section{Beneficios obtenidos}

Los principales beneficios son:

\begin{itemize}
  \item reproducibilidad del entorno y del caso de uso;
  \item separación entre metadatos técnicos y curación funcional;
  \item núcleo Python único para CLI, scripts y web;
  \item validación formal mediante SHACL;
  \item trazabilidad entre SQL, OpenMetadata, hoja funcional, JSON-LD e informe de validación;
  \item facilidad de operación mediante consola web.
\end{itemize}

\section{Limitaciones}

La solución tiene limitaciones deliberadas:

\begin{itemize}
  \item una tabla SQL no equivale semánticamente a un \texttt{dcat:Dataset}; en el caso de uso se usa como aproximación técnica gobernable;
  \item se genera una única \texttt{Distribution} y un único \texttt{DataService} por dataset;
  \item la clasificación HVD es una hipótesis de validación dentro de la plataforma, no una calificación jurídica automática;
  \item la validación se centra en metadatos, no en profiling de contenido ni calidad fila a fila;
  \item no se publica automáticamente en datos.gob.es, CKAN ni data.europa.eu;
  \item el modelo OpenMetadata se mantiene mínimo y no cubre todos los campos opcionales de DCAT-AP-ES.
\end{itemize}

La primera limitación merece una precisión conceptual. Un \texttt{dcat:Dataset} es una abstracción editorial: una colección de datos publicada y mantenida por un agente como una unidad coherente, con independencia de cómo se almacene físicamente~\cite{w3cDcat3}. Una tabla \ac{SQL}, en cambio, es una estructura física de almacenamiento. La correspondencia «una tabla es un dataset» es, por tanto, una decisión de modelado y no una equivalencia natural: a nivel técnico un dataset podría corresponder igualmente a una vista, a una vista materializada, al resultado de una consulta o a la unión lógica de varias tablas. Conceptos como un cursor o un índice no son candidatos válidos, porque un cursor es un mecanismo de recorrido y un índice una estructura de acceso, no un conjunto de datos publicable. En el caso de uso de validación se elige la tabla \texttt{gold} como unidad publicable por ser la representación más estable y autocontenida del dato refinado; en un escenario real, la hoja de gobierno permitiría asociar el dataset a la vista o agregación que mejor represente la unidad de publicación, sin alterar el núcleo.

Estas limitaciones no invalidan el trabajo: forman parte del alcance funcional declarado y delimitan la frontera entre lo que la plataforma demuestra de forma reproducible y lo que correspondería a una iniciativa de portal de datos abiertos a escala institucional, con presupuesto, equipo y compromisos de servicio diferenciados.

\section{Amenazas a la validez}

La principal amenaza a la validez externa es que el caso de uso trabaja con datos sintéticos y un número reducido de datasets publicables. La arquitectura, sin embargo, está diseñada para escalar a más tablas \texttt{gold} mediante la misma hoja funcional y el mismo workflow.

Conviene aclarar que la conformidad \ac{SHACL} «con warnings permitidos» no constituye una amenaza a la validez. En \ac{SHACL}, una violación con severidad \texttt{sh:Violation} indica el incumplimiento de una restricción obligatoria; un \texttt{sh:Warning}, en cambio, señala únicamente que el perfil \emph{recomienda} cubrir propiedades adicionales —recomendadas u opcionales— que no son obligatorias~\cite{w3cShacl}. El catálogo exportado no presenta violaciones bloqueantes: los warnings solo invitan a enriquecer metadatos no exigidos por el perfil. Aumentar esa cobertura mejora la calidad editorial del catálogo, pero su ausencia no invalida la conformidad obligatoria demostrada.

\section{Posicionamiento frente a entornos Big Data}\label{sec:posicionamiento-bigdata}

El enunciado oficial sitúa el trabajo en ``entornos Big Data y de computación en la nube'' \cite{uclmFichaTfe2026}, expresión que conviene matizar para evitar una lectura equívoca. La plataforma no opera en la capa de procesamiento masivo de datos, sino en la capa de gobierno de metadatos e interoperabilidad. El caso de uso de validación trabaja con un volumen reducido de datos sintéticos y, por tanto, no constituye Big Data en el sentido estricto de volumen, velocidad y variedad. Esta frontera es coherente con la distinción establecida en el capítulo~\ref{cap:objetivos}: el objeto gobernado son los metadatos, no los datos de negocio.

La relevancia para entornos Big Data es, en consecuencia, indirecta pero sustantiva. Los ecosistemas de datos a gran escala se caracterizan por activos heterogéneos, distribuidos y de múltiples fuentes, cuya utilidad depende de que permanezcan descubribles, descritos y reutilizables. Esa función —catalogación, descripción semántica e interoperabilidad— es precisamente la que DAMA encuadra como gestión de metadatos dentro del gobierno del dato \cite{damaDmbok}. La plataforma demuestra dicho mecanismo de forma reproducible: descubre activos técnicos, los describe, los exporta como catálogo DCAT-AP-ES y valida su conformidad.

En un escenario Big Data real, el proyecto se utilizaría sin cambios en su lógica de negocio, dado que el flujo opera sobre metadatos y es independiente del volumen del dato subyacente. El crecimiento se absorbería por tres vías: (i) la incorporación de nuevas fuentes y servicios en OpenMetadata, catálogo concebido para inventarios técnicos amplios y heterogéneos \cite{openmetadataDocs}; (ii) la ampliación del contrato funcional añadiendo más tablas \texttt{gold} a la misma hoja de gobierno y al mismo workflow idempotente, sin modificar el núcleo; y (iii) el despliegue nativo en la nube, que habilita el escalado de la capa de catálogo y gobierno conforme a las primitivas de orquestación de Kubernetes \cite{kubernetesConcepts}. La salida interoperable en DCAT-AP-ES es, además, el vehículo por el que un patrimonio de datos extenso se hace federable y localizable entre portales nacionales y europeos \cite{datosGobDcatApEs2025,w3cDcat3}. En síntesis, el trabajo no escala procesando más datos, sino gobernando más metadatos, que es la dimensión en la que el gobierno aporta valor a un entorno Big Data.

\section{Gobierno supervisado y responsabilidad del operador}\label{sec:gobierno-supervisado}

Aunque el flujo es ampliamente automatizable e idempotente, una parte esencial del gobierno reside en decisiones que deben permanecer bajo autoridad humana. Determinar qué activos son publicables, asignar la temática NTI-RISP \cite{datosGobNtiRisp}, clasificar un dataset como de alto valor o redactar su título, descripción y publicador no son transformaciones deterministas, sino juicios editoriales y de responsabilidad. El carácter manual y dinámico de esta curación no es una carencia del trabajo, sino una propiedad deseada del gobierno del dato.

El marco DAMA es explícito en este punto: el gobierno del dato asigna responsabilidad (\emph{accountability}) a roles definidos —propietario y \emph{steward} del dato— y la automatización debe asistir esa función, no sustituirla \cite{damaDmbok,datosGobDama2021}. Por consiguiente, la persona que opera o gobierna la plataforma debe conocer y aplicar las buenas prácticas del dominio, y responder de aquello que publica. La interoperabilidad técnica no exime de la responsabilidad sobre el contenido publicado.

Este principio tiene una implicación normativa concreta. La condición de conjunto de datos de alto valor está determinada por el Reglamento de Ejecución (UE) 2023/138 \cite{euHvdRegulation}; la plataforma trata la clasificación HVD como hipótesis de validación y no como una calificación jurídica automática, como ya se reconoció entre las limitaciones de este capítulo. Por ello, la categorización debe ser revisada y autorizada por una persona responsable antes de cualquier publicación efectiva, y nunca derivarse de forma ciega por mera configuración.

El diseño de la plataforma materializa esta filosofía mediante salvaguardas explícitas. La asistencia de autorrellenado solo completa campos vacíos con valores de validación y nunca sobrescribe la curación introducida por el operador, de modo que la automatización ayuda sin desplazar el criterio humano. Del mismo modo, las líneas de trabajo futuro orientadas a la ejecución periódica mediante un planificador o CI/CD, que se recogen a continuación, deben entenderse como automatización de la reproducibilidad, no como delegación de la decisión de publicar. La recomendación operativa es inequívoca: que el sistema sea técnicamente capaz de automatizar el flujo de extremo a extremo no implica que deba ejecutarse sin supervisión; la publicación de metadatos gobernados, por su visibilidad externa y por sus consecuencias jurídicas y reputacionales, debe ser un acto autorizado y supervisado.

\section{Trabajo futuro}

Las líneas de evolución más razonables son:

\begin{itemize}
  \item publicar automáticamente el JSON-LD en un CKAN o portal compatible cuando exista un destino real;
  \item ampliar la cobertura de metadatos recomendados y opcionales de DCAT-AP-ES;
  \item incorporar aprobaciones editoriales y roles si la web pasa a uso multiusuario;
  \item automatizar ejecuciones periódicas mediante Airflow, CI/CD o un scheduler;
  \item validar el comportamiento con más fuentes técnicas y más datasets \texttt{gold}, por ejemplo mediante ingesta desde APIs o bases de datos públicas.
\end{itemize}
